\documentclass[11pt,a4paper]{article}
\usepackage[T1]{fontenc}
\usepackage{isabelle,isabellesym}
\usepackage{amsmath,amssymb}
\usepackage{longtable}
\usepackage{pdfsetup}
\urlstyle{rm}
\isabellestyle{it}
\begin{document}
\title{The Combinatorial Four Colour Theorem}
\author{Ari Gilchrist}
\date{25 November 2026}
\maketitle
\begin{abstract}
We formalise the combinatorial Four Colour Theorem in Isabelle/HOL:
every finite planar bridgeless hypermap has a face colouring with four
colours. By duality, every finite planar loopless hypermap has a vertex
colouring with four colours. The development follows Gonthier's Rocq
formalisation of the Robertson--Sanders--Seymour--Thomas proof. It includes
the structural reduction, Kempe-chain arguments, reducibility of 633
configurations, and the discharging and unavoidability argument.
Executable certificate checkers are proved sound in HOL, and their finite
obligations are discharged using Isabelle's standard evaluation mechanism.
The final theorems have no additional reducibility or unavoidability
hypotheses. A dependency audit rejects skipped proofs and oracle dependencies
other than the standard code-evaluation oracle.
\end{abstract}
\tableofcontents
\parindent 0pt\parskip 0.5ex

\section{Statement and representation}
The Four Colour Theorem is represented by finite combinatorial maps,
following Gonthier~\cite{gonthier2008,gonthier-fourcolor}.
A hypermap consists of a finite carrier of darts and three permutations,
\texttt{edge}, \texttt{node}, and \texttt{face}, satisfying
\texttt{node (face (edge x)) = x}.
The permutations act as the identity outside the carrier.
Their orbits describe the edges, vertices, and faces. Planarity is expressed
by genus zero, with genus defined from the orbit counts and the Euler
characteristic. The proved Euler and Jordan results justify the structural
use of this definition.

The datatype \texttt{color} has exactly four constructors.
A face colouring is constant along each face permutation and changes
colour along an edge permutation, for every dart in the carrier.
The bridgeless assumption excludes an edge joining a face to itself.
For the dual graph formulation, the corresponding assumption is looplessness.
The main results in theory \texttt{Four\_Colour} are:
\begin{itemize}
\item \texttt{four\_color\_hypermap}: a planar bridgeless hypermap is
  four-colourable;
\item \texttt{hypermap.four\_colour\_theorem}: the same result with the
  hypermap, planarity, and bridgelessness assumptions displayed separately;
\item \texttt{hypermap.four\_colour\_graph}: a planar loopless hypermap
  admits a vertex colouring with four colours.
\end{itemize}
The result is the combinatorial map formulation. No geometric embedding
predicate or theorem about arbitrary subsets of the real plane is asserted
by this entry.

\section{Proof structure}
The finite-list, path, and orbit libraries follow the conventions of
SSReflect and Mathematical Components~\cite{mathcomp}.
The mathematical argument is the reducibility and discharging proof of
Robertson, Sanders, Seymour, and Thomas~\cite{rsst1997}, as formalised by
Gonthier. The port preserves the 633 configuration programs and the
presentation certificates of that development. Checkers validate these
programs and certificates in Isabelle; correctness does not depend on the
program that translated the data.

\begin{longtable}{p{0.28\textwidth}p{0.64\textwidth}}
\textbf{Session} & \textbf{Contents} \\
\hline
\texttt{Four\_Colour\_Base} & Finite permutations, hypermaps, Euler and Jordan
properties, colourings, cubic reduction, and patches. \\
\texttt{Four\_Colour\_Maps} & Kempe arguments, chromogram trees, configuration
maps, contractions, reducibility, embeddings, and quizzes. \\
\texttt{Four\_Colour\_Parts} & Parts, reducibility of parts, symmetries,
discharging, and hubcaps. \\
\texttt{Four\_Colour\_Presentations} & Configuration data and the presentation
checker. \\
\texttt{Four\_Colour\_Checkers} & Soundness of executable checkers and the
logical final reductions. \\
\texttt{Four\_Colour\_Runtime} & Verified numeric code equations and shared
generated checker code. \\
\texttt{Four\_Colour\_Quizzes} & Quiz certificates for all 633 configurations. \\
\texttt{Four\_Colour\_Reducibility} & Reducibility certificates for all
633 configurations. \\
\texttt{Four\_Colour\_Unavoidability} & Score caps and presentations for
arities 5 through 11. \\
\texttt{Four\_Colour} & Unavoidability, the final map and graph theorems,
and the dependency audit. \\
\end{longtable}
The session boundaries allow the certificate computations to run in separate
processes and to be retained in intermediate heaps. This is important on
machines with limited memory. They are packaging boundaries, not additional
mathematical assumptions.

\section{Computation and trust}
The soundness lemmas for the certificate checkers are proved in HOL. The
individual concrete obligations use \texttt{eval}. Its oracle dependency is
\texttt{Code\_Generator.holds\_by\_evaluation}, so evaluation additionally
trusts Isabelle's code generator and the ML runtime. This is not an
oracle-free kernel reconstruction of all certificate evaluations.
\texttt{code\_reflect} shares generated code to avoid recompiling the
checkers for every configuration. Numeric code equations are proved before
use; no external executable provides assumed certificate results.

The final ML audit checks the actual dependency graphs of the public theorems.
It rejects skipped proofs and any oracle other than the standard evaluation
oracle. No configuration checker, presentation certificate, or absence of
minimal counterexamples remains as an assumed premise of the final theorems.

\section{Sources and authorship}
The development is an Isabelle/HOL port of the upstream \texttt{fourcolor}
repository at revision
\texttt{c1d6b1cd5288bea4b067aac13cdde3c18dffe018}.
The upstream mathematical work and Rocq formalisation are due to Georges
Gonthier and the original contributors; the original proof is due to
Robertson, Sanders, Seymour, and Thomas. This entry does not claim a new
mathematical proof strategy.

The Isabelle port was developed using Anthropic Claude and OpenAI Codex.
Claude produced much of the initial formalisation. Codex contributed earlier
sewing and snipping proofs, then completed remaining proofs, computational
certificate checks, final assembly, and submission preparation after the
recorded handoff. The repository's authorship record distinguishes these
contributions. The upstream copyright and licence notices are preserved in
\texttt{NOTICE} and \texttt{LICENSES/CeCILL-B.txt}; new material is covered as
specified in \texttt{LICENSE}.

\section{Formal development}
Long generated configuration and presentation literals are omitted from this
PDF; their complete source remains in the entry. The following sections
include the supporting theories in dependency order and the final assembly.
\input{FC_Seq}
\input{FC_Fingraph}
\input{FC_Perm}
\input{Hypermap}
\input{Walkup}
\input{Jordan}
\input{Euler_Tree}
\input{Geometry}
\input{Color}
\input{Coloring}
\input{Cube}
\input{Mincex}
\input{Patch}
\input{Chromogram}
\input{Gtree}
\input{Ctree}
\input{Initctree}
\input{Initgtree}
\input{Ctreerestrict}
\input{Gtreerestrict}
\input{Sew}
\input{Snip}
\input{Kempe}
\input{Revsnip}
\input{Cfmap}
\input{Cfmap_ext}
\input{Cfmap_ka}
\input{Cfmap_prog}
\input{Cfcolor}
\input{Kempetree}
\input{Hypermap_iso}
\input{Birkhoff_check}
\input{Birkhoff}
\input{Spoke}
\input{Cfmap_ker}
\input{Cfcontract}
\input{Cfcontract2}
\input{Cfreducible}
\input{Contract}
\input{Contract2}
\input{Embed_pre}
\input{Embed}
\input{Embed1}
\input{Embed2}
\input{Quiz}
\input{Cfquiz}
\input{Quiztree}
\input{Part}
\input{Redpart}
\input{Redpart2}
\input{Redpart3}
\input{Converse}
\input{Mirror}
\input{Discharge}
\input{Hubcap}
\input{Configurations}
\input{Present}
\input{PresentBase}
\input{Check_Map}
\input{Quiz_Embed}
\input{Four_Colour_Reduction}
\input{Unavoid}
\input{Certificate_Runtime}
\input{Quiz_Certificates}
\input{Reducibility}
\input{Score_Caps}
\input{Present5}
\input{Present6}
\input{Present7}
\input{Present8}
\input{Present9}
\input{Present10}
\input{Present11}
\input{Combinatorial_Four_Colour}
\input{Realplane}
\input{Colour_Compactness}
\input{Finitize}
\input{Grid_Geometry}
\input{Grid_Rectangles}
\input{Approx}
\input{Map_Representatives}
\input{Matte}
\input{Matte_Extension}
\input{Scaled_Matte}
\input{Adjacency_Boxes}
\input{Matte_Extension2}
\input{Coarse_Geometry}
\input{Matte_Contact}
\input{Coarse_Extension}
\input{Matte_Order}
\input{Matte_Refinement}
\input{Connected_Matte}
\input{Region_Mattes}
\input{Plane_Discretization_Data}
\input{Matte_Port_Audit}
\input{Grid_Adjacency}
\input{Grid_Bounds}
\input{Grid_Contact}
\input{Grid_Permutations}
\input{Grid_Rotations}
\input{Framed_Grid}
\input{Grid_Counting}
\input{Grid_Vertices}
\input{Grid_Darts}
\input{Grid_Connectivity}
\input{Grid_Planarity}
\input{Grid_Rings}
\input{Grid_Disks}
\input{Grid_Cutout}
\input{Grid_Colouring}

\input{session}
\bibliographystyle{abbrv}
\bibliography{root}
\end{document}
